\section{Postprocessing}
\label{sec:alignment-postprocessing}

After the alignment is done, BioseqDB computes CIGAR string, extracts all values returned by the final aligner (like exact coordinates, score etc.), and also calculates
some additional values (if they are requested by the user in query):
\begin{itemize}
  \item \textbf{Divergence} Matching minimizers give an estimate of the fraction of differing bases
  \item \textbf{Mapping quality (MAPQ)} This estimates how subjectively good are the mapping locations (are we confident enough that this is a match)
\end{itemize}

Minimap defines diverges exactly as mentioned before in Section~\ref{sec:minimizer-divergence-history}:
\begin{equation}
  \widehat{\epsilon}=\frac{1}{k}\ln\frac{n}{m},
  \qquad
  \widehat{d}=1-\exp(-\widehat{\epsilon})
            =1-\left(\frac{m}{n}\right)^{1/k}.
  \label{eq:postprocessing-divergence}
\end{equation}
where $n$ is the number of minimizers that are part of chains and $m$ is the number of all query minimizers between first match and last match
BioseqDB similarily to what minimap2 does, returns
\(\widehat{d}\), which is approximation of \(\widehat{\epsilon}\) assuming low sequence divergence.
To obtain \(n\) in BioseqDB I uise cumulative counts of minimizers by position, so subtracting two counts gives the number in an interval without
sorting the emitted minimizers (this avoids cirtical problem I found earlier, which made this step thousands of time slower - Section~\ref{sec:minimizer-divergence-history}).

For MAPQ, if two reference locations give nearly equal alignment scores, BioseqDB assigns
low MAPQ, because both alignments may match well and we cannot be confident enough that this location is the correct one.

The maximum MAPQ in minimap2 is 60 and it's similar to usual phred score scale used in base calling (for example in SAM formats), where $P=10^{-Q/10}$~\autocite{sam_format_specification}
equal to 60 means 1 in million chance for error.

BioseqDB uses a direct port of minimap2's heuristic~\autocite{minimap2_mapq_implementation}.
The equation is given below. Please note that alignment score is different to score given by chaining stage:
\begin{equation}
  \begin{aligned}
    \mathrm{MAPQ}\approx\operatorname{clip}_{C_{\min},C_{\max}}\Bigl(
      &\min\!\left\{
        C_1Ip(1-\rho^{C_{\mathrm{power}}})\ln\frac{D_{\mathrm{pri}}}{a},\,
        C_2I^{C_{\mathrm{power}}}\frac{D_{\mathrm{pri}}-D_{\mathrm{sec}}}{a}
      \right\}\\
      &{}-C_3\ln(N+1)
    \Bigr).
  \end{aligned}
  \label{eq:mapq-final}
\end{equation}
\begin{itemize}
  \item \(D_{\mathrm{pri}}\): best DP score returned by the aligner for the
    primary chain.
  \item \(D_{\mathrm{sec}}\): highest DP score returned by the aligner among
    secondary chains competing with the primary chain for an overlapping
    region of the query.
  \item \(a\): reward for one matching base.
  \item \(I\): factor of how many bases we matched (matching bases divided by alignment block length)
  \item \(p\): MAPQ multiplier factor calculated as $p=\operatorname{clip}_{1,\infty}(\frac{s_{\pathrm{pri}}}{100}, \frac{{N}_{\pathrm{pri}}}{10})$, where ${s}_{\pathrm{pri}}$ is the chaining score of primary chain and ${N}_{\pathrm{pri}}$ number of anchors in that chain.
  It's used to penalize short or repetitive chains.
  \item \(\rho=(D_{\mathrm{sec}}/D_{\mathrm{pri}})
    \max(s_{\mathrm{sec}},s_{\min})/s_{\mathrm{orig}}\): ratio combining alignment and chaining scores. 
    Here \(s_{\mathrm{sec}}\) is the best secondary chaining score, \(s_{\min}\) is the configurable fixed parameter used by floor for the secondary chaining score and \(s_{\mathrm{pri}}\) is
    the primary chaining score.
  \item \(N\): number of secondary chains for that fragment.
  \item \(\operatorname{clip}_{C_{\min},C_{\max}}\) shorthand notation for $min(max(x, C_{\min}), C_{\max})$
  \item BisoeqDB defaults to minimap2's defaults: \(C_1=40\), \(C_2=6.02\), \(C_3=4.343\),
    \(C_{\mathrm{power}}=2\), \(C_{\min}=0\), \(C_{\max}=60\).
\end{itemize}

\info{MAPQ implementation details}{
  In minimap2 it's possible to disable base alignment and when that happens certain adjustments happen, for example $D_{\mathrm{sec}}/D_{\mathrm{pri}}$ becomes one,
  instead of $D$ we use $s$ where possible.
  The origin of the formula is empirical and unfortunately I didn't go into details understanding completely its structure. 
  Some details can be found in original minimap2 paper~\autocite{minimap2_paper}.
}
